\chapter{Related Work} \label{sec:related}
Es gibt viele Studien, welche die Echtzeitfähigkeit von Wi-Fi untersuchen.
Angefangen mit Sudien, welche sich mit Wi-Fi und harter Echtzeit auseindaer setzen.
Dabei setzen sowohl Wei et al. \cite{6728869}, als auch Moraes et al. \cite{4638758} auf das \ac{TDMA} Verfahren.
In beiden Studien wird berichtet, dass harte Echtzeit nur mittels einer deterministischer \ac{MAC} erreicht wird.
Somit kann in \cite{6728869} eine Frequenz an Echtzeitpaketen von 6Khz erreicht werden.

Avdotin et al. setzt in \cite{Avdotin_2019} ebenfalls an der \ac{MAC} an, jedoch spezifisch für \ac{OFDMA}.
In dieser Studie wird die Resource Allocation von \ac{OFDMA} angepasst und harte Echtzeitfähikgeit kann bis zu eine Zykluszeit von 1ms gehalten werden.
Dies wiedersprcith den Aussagen von \cite{6728869} und \cite{4638758}, dass nur mittels \ac{TDMA} harte Echtzeit erreicht werden kann.
Aber dennoch ist ein Eingriff in die Firmware nötig.

Auch zu den Auswirkungen von \ac{QoS} Parametern auf die Echtzeitfähigkeit existieren Studien.
Somit erörtern Tsolakou et al. in \cite{Tsolakou_2002s} die Auswirkungen von verschiedenen \ac{Qdisc} Disziplinen auf die Echtzeitfähigkeit.
Des Weiteren legt Dahoud et al. in \cite{da_2019} die Auswirkungen von verschiedenen \ac{DSCP}-Werten auf die Latenz dar.

Aber auch zur Allmeinen Performance existieren Studien.
Somit beschreit Chotalia et al. in \cite{10.1007/978-3-031-40564-8_14} allgemeine Performance metrikend er Standarts 4 - 6 .
Dabei geht er besonders auf den erzeugbaren Datendurchsatz ein.

Dabei stellt sich nun die Frage, wie diese Arbeiten und Erkenntnisse von den Ergebenissen in dieser Arbeit zu separieren sind.
Das Ziel in dieser Arbeit sit es, wie in Kapitel \ref{subsec:objective} beschrieben, die minimale Zyklusszeit für Echtzeitdaten herauszufinden, ohne dabei den Wi-Fi Standard grundlegend zu ändern.
Somit sind die Ansätze von \ac{TDMA} und der Resoure Allocation Änderung von \ac{OFDMA} in dieser Arbeit nicht enthalten, da diese einen Eingriff in die Firmware vorraussetzwen würden.
Des weiteren sind viele Studien, wie \cite{10.1007/978-3-031-40564-8_14} nur in Simulationen oder reinen Laborumgebungen enstanden.
In dieser Arbeit sollen ebenfalls realisitsche Umgebungen und deren Einfluss auf die Datenverbindung betrachtet werden.
Abschließend sollen nicht nur Echtzeitdaten übertragen, sondern auch die gleichzeitige Nutzung normaler Daten maximiert werden.
Die KOmbination aus Betrieb in LAbor und Laborfreiem raum, keine Änderung am Standard und Maximierung von gleichzeitigen Nutzdaten, in KOmbination von Tests von verschiedenen Standards und \ac{QoS} Kombinationen, wurde bisher noch nicht untersucht.

Im Folgenden wird darauf eingegangen, wie der Testaufbau in dieser Arbeit umgesetzt ist.
